% from ega1-9-fr.tex
\oldpage[I]{171}
\[
  \xymatrix{
    X_{(S')}\ar[r]^{f_{(S')}}\ar[d] &
    Y_{(S')}\ar[d]\\
    X\ar[r]_f &
    Y
  }
\]
\end{proof}

\begin{corollary}[9.1.11]
\label{I.9.1.11-fr}
Soient $X$, $Y$ deux $S$-préschémas, $\mathscr{F}$ (resp. $\mathscr{G}$) un
$\mathscr{O}_X$-Module (resp. $\mathscr{O}_Y$-Module) quasi-cohérent.
L'image réciproque du faisceau
$(\mathscr{F}_{(S')}\otimes_{S'}\mathscr{G}_{(S')})$ par l'isomorphisme
canonique
$(X\times_S Y)_{(S')}\xrightarrow{\sim}
(X_{(S')})\times_{S'}(Y_{(S')})$
(\hyperref[I.3.3.10-fr]{3.3.10}) est égale à
$(\mathscr{F}\otimes_S\mathscr{G})_{(S')}$.
\end{corollary}

\begin{proof}
Si $p$, $q$ sont les projections de $X\times_S Y$, l'isomorphisme en
question n'est autre que $(p_{(S')},q_{(S')})_{S'}$~; le corollaire résulte
des prop. (\hyperref[I.9.1.4-fr]{9.1.4}) et
(\hyperref[I.9.1.10-fr]{9.1.10}).
\end{proof}

\begin{proposition}[9.1.12]
\label{I.9.1.12-fr}
Avec les notations de (\hyperref[I.9.1.2-fr]{9.1.2}) soient $z$ un point de
$X\times_S Y$, $x=p(z)$, $y=q(z)$~; la fibre
$(\mathscr{F}\otimes_S\mathscr{G})_z$ est isomorphe à
$(\mathscr{F}_x\otimes_{\mathscr{O}_x}\mathscr{O}_z)
\otimes_{\mathscr{O}_z}
(\mathscr{G}_y\otimes_{\mathscr{O}_y}\mathscr{O}_z)
=\mathscr{F}_x\otimes_{\mathscr{O}_x}\mathscr{O}_z
\otimes_{\mathscr{O}_y}\mathscr{G}_y$.
\end{proposition}

\begin{proof}
Comme on peut se ramener au cas affine, la proposition résulte de la formule
(\hyperref[I.1.6.5.1-fr]{1.6.5.1}).
\end{proof}

\begin{corollary}[9.1.13]
\label{I.9.1.13-fr}
Si $\mathscr{F}$ et $\mathscr{G}$ sont de type fini, on a
\[
  \operatorname{Supp}(\mathscr{F}\otimes_S\mathscr{G})
  =p^{-1}(\operatorname{Supp}(\mathscr{F}))
  \cap q^{-1}(\operatorname{Supp}(\mathscr{G})).
\]
\end{corollary}

\begin{proof}
Comme $p^*(\mathscr{F})$ et $q^*(\mathscr{G})$ sont de type fini sur
$\mathscr{O}_{X\times_S Y}$, on est ramené, par
(\hyperref[I.9.1.12-fr]{9.1.12}) et
(\hyperref[0.1.7.5-fr]{0, 1.7.5}), au cas particulier où
$\mathscr{G}=\mathscr{O}_Y$, c'est-à-dire à démontrer la formule
\[
  \operatorname{Supp}(p^{-1}(\mathscr{F}))
  =p^{-1}(\operatorname{Supp}(\mathscr{F})).
  \tag{9.1.13.1}\label{I.9.1.13.1-fr}
\]

Le même raisonnement que dans (\hyperref[0.1.7.5-fr]{0, 1.7.5}) ramène à
vérifier que l'on a, pour tout $z\in X\times_S Y$,
$\mathscr{O}_z/\mathfrak{m}_x\mathscr{O}_z\neq0$ (avec $x=p(z)$), ce qui
découle du fait que l'homomorphisme
$\mathscr{O}_x\to\mathscr{O}_z$ est \emph{local} par hypothèse.
\end{proof}

Nous laissons au lecteur le soin d'étendre à un produit d'un nombre
quelconque de facteurs les résultats démontrés dans ce numéro pour deux
facteurs.

\subsection{Image directe d'un faisceau quasi-cohérent.}
\label{subsection:I.9.2-fr}

\begin{proposition}[9.2.1]
\label{I.9.2.1-fr}
Soit $f:X\to Y$ un morphisme de préschémas. On suppose qu'il existe un
recouvrement $(Y_\alpha)$ de $Y$ par des ouverts affines ayant la propriété
suivante~: chacun des $f^{-1}(Y_\alpha)$ admet un recouvrement fini
$(X_{\alpha i})$ par des ouverts affines contenus dans
$f^{-1}(Y_\alpha)$, tel que chacune des intersections
$X_{\alpha i}\cap X_{\alpha j}$ soit elle-même réunion finie d'ouverts
affines. Dans ces conditions, pour tout $\mathscr{O}_X$-Module
quasi-cohérent $\mathscr{F}$, $f_*(\mathscr{F})$ est un
$\mathscr{O}_Y$-Module quasi-cohérent.
\end{proposition}

\begin{proof}
La question étant locale sur $Y$, on peut supposer $Y$ égal à l'un des
$Y_\alpha$, donc supprimer les indices $\alpha$.

\begin{enumerate}
\item[a)] Supposons d'abord que les $X_i\cap X_j$ soient eux-mêmes des
ouverts \emph{affines}. Posons
$\mathscr{F}_i=\mathscr{F}|X_i$,
$\mathscr{F}_{ij}=\mathscr{F}|(X_i\cap X_j)$, et soient
$\mathscr{F}'_i$ et $\mathscr{F}'_{ij}$ les images de $\mathscr{F}_i$ et
$\mathscr{F}_{ij}$ respectivement par les restrictions de $f$ à $X_i$ et
$X_i\cap X_j$~; on sait que les $\mathscr{F}'_i$ et
$\mathscr{F}'_{ij}$ sont quasi-cohérents
(\hyperref[I.1.6.3-fr]{1.6.3}). Posons
$\mathscr{G}=\bigoplus_i\mathscr{F}'_i$,
$\mathscr{H}=\bigoplus_{i,j}\mathscr{F}'_{ij}$~; $\mathscr{G}$ et
$\mathscr{H}$ sont des $\mathscr{O}_Y$-Modules quasi-cohérents~; nous allons
définir un homomorphisme $u:\mathscr{G}\to\mathscr{H}$ tel que
$f_*(\mathscr{F})$ soit le \emph{noyau} de $u$~; il en résultera que
$f_*(\mathscr{F})$ est quasi-cohérent
(\hyperref[I.1.3.9-fr]{1.3.9}). Il suffit de définir $u$
